\documentclass[border=5pt]{standalone}
% 命名色表必须在 pgfplots 之前声明：pgfplots 内部会先加载 xcolor，
% 之后再 \usepackage[dvipsnames,svgnames]{xcolor} 会触发 Option clash
\PassOptionsToPackage{dvipsnames,svgnames}{xcolor}
\usepackage{pgfplots}
\usepackage{pgf-pie} % 饼图：AI 代码可直接使用 \pie
\pgfplotsset{compat=1.18}
\usepgfplotslibrary{fillbetween}  % 面积图 / 堆叠面积图 / \closedcycle 填充路径
% 注：error bars 是 pgfplots 内置功能（在核心 pgfplots.errorbars.code.tex），不需要单独 \usepgfplotslibrary 加载
\usepackage{amsmath}
\usepackage{amssymb}

% 支持中文
\usepackage{fontspec}
\usepackage{xeCJK}
\setCJKmainfont{SimSun}
\setmainfont{Times New Roman}

\begin{document}

\begin{tikzpicture}
\pie[radius=3, text=legend, style={font=\small}, text width=2cm]{32/品牌A, 26/品牌B, 18/品牌C, 14/品牌D, 10/品牌E}
\node[anchor=north west, font=\scriptsize] at (axis description cs:0.0,-0.15) {数据来源：03-pie-share.xlsx};
\end{tikzpicture}

\end{document}